\scanpage{28}
Alors la suite spectrale implique
\[
H^i_Y(X,\mu_{nX}^{\otimes d})=0\quad\text{si }i<2d
\]
\[
H^{2d}_Y(X,\mu_{nX}^{\otimes d})\simeq
\Gamma(Y,\underline H^{2d}_Y(\deleted{\unsure{\mathbb Z}}\mu_{nX}^{\otimes d}))
\]
d'où un \deleted{hom canonique} élément canonique
\[
\unsure{\gamma|\sigma}^{(n)}_{X/Y}\in H^{2d}_Y(X,\mu_{nX}^{\otimes d})
\]
[qu'on peut maintenant spécialiser dans
$H^{2d}_{\deleted{\unsure{Y}}}(X,\mu_{nX}^{\otimes d})$, ou si
$X$ est propre sur un $S$, dans
\[
H^{2d}_{\Phi_{Y/X}}(X,\mu_{nX}^{\otimes d}).
\]
\margin{$\Phi_{Y/X}\longrightarrow$ famille de spts}
\margin{N.B. \unsure{si} $(\beta_{Y/X}\in H^1_Y(X,G_{mX})$
et $\unsure{\gamma|\sigma}^{(n)}_{Y/X}=\partial_n(\beta_{Y/X})$,
en particulier son image dans $H^2_Y(X,G_{mX})$ est nulle}

\uline{Exemple 1}. $X$ lisse sur $k$ \inserted{sép. clos.},
\deleted{$x$ pt} dim relative $d$,
\deleted{\unsure{$\forall x$ pt de $X$} \inserted{\unsure{pt sur $k$, d'où}}
\unsure{on obt.|on obtient} \unsure{des C fl. de}}
\[
\unsure{\gamma|\sigma}_{x/X}\in H^{2d}_!(X,\mu_{nX}^{\otimes d}).
\]
Il va résulter du théorème qui va suivre que \deleted{pour $\unsure{\forall|\text{tt}}$}
si $X$ est géom. connexe, cet élément est
\uline{une base de} $H^{2d}_!(X,\mu_{nX}^{\otimes d})$
\uline{et ne dépend pas du choix de $x$};
\deleted{\unsure{ce}} C'est qu'on avait vérifié directement dans le cas où
$d=1$, en \deleted{\unsure{le cas où}} \inserted{nous ramenant}
au cas où $C$ étant complète et où on dispose déjà
\mspage{27}
